\scanpage{32}%
\mspage{17}%
and the law of composition by
\[
(u',s',t')\mathbin{\handwritten{o}}(u,s,t)
=(u\,u',u's+t',u^{\raisebox{-.38ex}{$\scriptstyle\prime$}\,3}t+t'+u^{\raisebox{-.38ex}{$\scriptstyle\prime$}\,2}rs')\, .
\]

When the curve is put in form (6.8), these equations reduce to
\[
\begin{aligned}
u\in\mu_3 \qquad s\in\mathbb{F}_4 \qquad t^2+t=s^3
 &\quad ({}=0\ \text{or}\ 1\ \text{according as}\\
 &\qquad s=0\ \text{or}\ s\ne 0) \, .
\end{aligned}
\]

7. \quad \uline{Ring of integral modular forms}.

By an integral modular form of weight $n$ one means a rule which to each curve of genus one over a base $S$ assigns a section of $\omega^{\otimes n}$, and this compatibly with base change. It would amount\margin{\handwritten{\tikz[baseline=-0.2ex]{\draw[thin,gray] (0,-0.35em) -- (0.48em,1.30em);} \uuline{why} ?}} to the same thing to restrict to curves without additive fibres.

\uline{Proposition} 7.1. \quad \uline{The ring of integral modular forms is generated over} $\mathbb{Z}$ \uline{by} $c_4$, $c_6$ \uline{and} $\Delta$ \uline{subject to the single relation}
\[
c_4^3-c_6^2=1728\,\Delta \, . \tag*{(7.2.)}
\]
\margin{\tikz[baseline=(note.north)]{\node[inner sep=2pt,text width=9.8em,align=left] (note) {\handwritten{N.B. On the modular site, the divisors of $c_4$, $c_6$, $\Delta$ have \unsure{multiplicity|multiplicities} $1$ ; That of $\Delta$ is smooth over $\mathbb{Z}$, those of $c_4$, $c_6$ are so except above the points $2$ and $3$.}};\handwritten{\draw[thin,rounded corners=1pt] (note.north west) -- (note.south west) -- (note.south east);}}}

Applying the definition to the curve (1.3), one sees that an integral modular form of weight $n$ is expressed as a poly\tikz[baseline=(nome.base)]{\node[inner sep=0pt] (nome) {nomial};\handwritten{\draw[thin,gray,rotate around={-25:(nome.center)}] (nome.center) ellipse (1.04em and 1.65em);}} of weight $n$ in the $a_i$. With $2$ and $3$ invertible, it follows from 2.5 that these forms are \tikz[baseline=(des.base)]{\node[inner sep=0pt] (des) {some};\handwritten{\draw[thin,gray,rotate around={-15:(des.center)}] (des.center) ellipse (1em and 1.02em);}} poly\tikz[baseline=(noms.base)]{\node[inner sep=0pt] (noms) {nomial};\handwritten{\draw[thin,gray,rotate around={-25:(noms.center)}] (noms.center) ellipse (0.84em and 1.47em);}}s in $c_4$ and $c_6$, and in particular that $c_4$, $c_6$ and $\Delta$ satisfy no relations other than (7.2). It remains to prove that if a poly\tikz[baseline=(nomt.base)]{\node[inner sep=0pt] (nomt) {nomial};\handwritten{\draw[thin,gray,rotate around={-25:(nomt.center)}] (nomt.center) ellipse (0.85em and 1.53em);}} in $c_4$, $c_6$ and $\Delta$ is identically zero in characteristic $2$ (resp. $3$), then it is divisible by $2$ (resp.$3$) in $\mathbb{Z}[c_4,c_6,\Delta]/\handwritten{\bigl(}c_4^3-c_6^2-1728\,\Delta\handwritten{\bigr)}$, which means%
% archive sheet number: 31
